\chapter{Conclusiones}
\label{chap:conclusiones}

Este trabajo partió de la pregunta de si la facilitación humana de una discusión puede describirse con suficiente claridad para implementarla como un sistema. La revisión de la literatura permitió dividir el proceso en decisiones observables. A partir de ellas, se implementó un sistema que recibe un hilo de discusión, analiza su estado, decide si corresponde intervenir y genera una respuesta candidata cuando la intervención está justificada.

\section{Cumplimiento de los objetivos}

El objetivo principal era diseñar e implementar un sistema basado en LLMs y técnicas de ingeniería del contexto para analizar y facilitar discusiones académicas asíncronas procedentes de distintas fuentes. Este objetivo se cumplió en lo relativo al diseño y la implementación. El sistema utiliza el contenido del hilo de discusión, los objetivos de aprendizaje cuando se proporcionan y criterios derivados de la literatura para decidir si intervenir y producir una respuesta candidata. El contrato de entrada permite incorporar otra fuente mediante un adaptador, aunque esta posibilidad no se verificó con un segundo LMS o foro.

Los experimentos aportan evidencia de viabilidad técnica sobre ese objetivo porque la secuencia de facilitación se ejecutó con modelos locales y externos. Además, como las contribuciones se producen en momentos distintos, el sistema puede realizar el análisis en segundo plano y la ejecución no requiere una respuesta conversacional en tiempo real.

Los experimentos documentan la ejecución técnica del sistema. Sin embargo, la evaluación no incluyó estudiantes, instructores ni cursos en funcionamiento, por lo que no permite establecer efectos sobre el aprendizaje, la carga docente o la capacidad de atender más discusiones.

El primer objetivo específico proponía un sistema de agentes organizado en etapas. Cada etapa concentra una responsabilidad y conserva una salida que puede inspeccionarse. La separación permite modificar, por ejemplo, el comportamiento de un rol sin afectar la clasificación. Además, las interfaces de fuentes, proveedores de modelos y almacenamiento permiten cambiar esos servicios sin modificar los nodos. El requisito relativo a tiempos y consumo de recursos razonables no quedó verificado, ya que no se calculó el coste completo de cada intervención ni se estudió la distribución de su latencia.

No se ejecutó una comparación con una única llamada al mismo modelo, de modo que la evidencia disponible no permite afirmar que la arquitectura multiagente produzca mejores intervenciones que una llamada directa. La separación en agentes aporta responsabilidades explícitas, decisiones intermedias que pueden inspeccionarse y la posibilidad de modificar cada etapa por separado.

El segundo objetivo fue definir una estrategia de ingeniería del contexto fundamentada en la literatura. La definición de roles constituyó el punto de partida del diseño. Ante un proceso que inicialmente no era fácil de describir, estudiar cómo interviene un facilitador humano permitió convertir conocimientos pedagógicos en responsabilidades, criterios y técnicas. Esa traducción sirvió de base para la implementación, aunque sigue siendo una primera aproximación realizada desde la ingeniería. Para convertirla en una estrategia pedagógica validada será necesario incorporar a especialistas en educación y el conocimiento empírico acumulado en el aula.

El tercer objetivo, relativo a la independencia de la plataforma, se abordó como una demostración técnica. El contrato común se utilizó con escenarios reproducibles, seis hilos históricos anonimizados y el adaptador de Open edX. La prueba de concepto implementó la adaptación de los datos y el mecanismo de activación. La secuencia de nodos permanece fija en la implementación actual. También quedan pendientes la integración en el trabajo cotidiano del instructor y la publicación segura en un curso.

El cuarto objetivo fue evaluar el comportamiento del sistema. Las dos series de ejecuciones produjeron 46 intervenciones evaluables, que un LLM de propósito general valoró con una media de 3,911 sobre 5. El evaluador marcó 29 casos para revisión, 27 de ellos por información no sustentada en el hilo de discusión. Estas puntuaciones y alertas describen las salidas obtenidas, aunque no sustituyen el criterio de especialistas ni permiten concluir que una intervención sea pedagógicamente efectiva.

\section{Aprendizajes del diseño y la implementación}

Los seis hitos descritos en la introducción se completaron con distinto grado de alcance. La definición del problema y la revisión de la literatura permitieron formular el proceso, mientras que la implementación produjo el motor y el panel. El contrato común y Open edX permitieron trabajar con las fuentes, y las series de ejecuciones en Idril y AWS EC2 proporcionaron la evaluación técnica. Además, ambos entornos quedaron desplegados para la experimentación y la demostración.

La parte más difícil fue caracterizar el comportamiento de los modelos y construir el sistema alrededor de sus respuestas. Las configuraciones evaluadas combinaban diferencias en los modelos, el modo de solicitar las salidas estructuradas y los umbrales, por lo que los resultados no permiten atribuir las diferencias a un único factor. Los registros muestran fallos de ejecución y resultados parciales. En consecuencia, la implementación necesitó conservar las trazas, limitar los tiempos de ejecución, mantener los resultados parciales y adaptar el tratamiento de las salidas.

Los dos entornos permitieron estudiar condiciones diferentes. Los modelos externos ofrecieron menores tiempos medios de ejecución en estas pruebas, aunque las diferencias de infraestructura, proveedor y carga impiden una comparación directa. Idril proporcionó una alternativa basada en infraestructura propia. El coste completo de las intervenciones no se calculó, por lo que no es posible determinar qué opción sería más sostenible para una universidad. Los datos muestran que la misma arquitectura puede ejecutarse con modelos locales y externos.

\section{Contribución}

Este trabajo aporta tres elementos técnicos. El primero es una descomposición operativa de la facilitación de discusiones académicas asíncronas. A partir de los roles y criterios identificados en la literatura, el sistema separa la caracterización del hilo de discusión, la decisión de intervenir, la selección de un rol y una técnica y la generación de una respuesta candidata.

El segundo elemento es la implementación de esa secuencia mediante etapas que pueden inspeccionarse. Cada etapa conserva una salida estructurada con su razonamiento, y los mecanismos de validación y reintento mantienen los resultados parciales cuando una etapa posterior falla. Las fuentes de hilos de discusión, los proveedores de modelos y el almacenamiento se conectan mediante interfaces.

El tercer elemento es el mecanismo de trazabilidad utilizado para examinar el sistema. Cada archivo de resultados relaciona el hilo de discusión y la configuración con las decisiones intermedias, la respuesta, la duración y los errores registrados. El panel y los registros técnicos permiten consultar esa información, y la misma rúbrica se aplicó a las 46 intervenciones evaluables.

\section{Limitaciones}

La evaluación caracteriza el comportamiento técnico del sistema en las configuraciones estudiadas. Sus condiciones delimitan el alcance actual de los resultados y señalan las comprobaciones necesarias para continuar el trabajo.

\begin{itemize}
	\item El control de calidad cubrió dos series de ejecuciones, ocho configuraciones de modelo y 18 hilos de discusión. Esta primera cobertura permitió observar el funcionamiento de la secuencia y los fallos registrados. La evaluación de utilidad con estudiantes, instructores o cursos en funcionamiento queda pendiente.
	\item Los 18 hilos estaban redactados en inglés y la muestra no se diseñó para comparar contextos culturales ni normas de participación diferentes. La evaluación puede ampliarse con discusiones en español, variantes lingüísticas y comunidades educativas con otras convenciones.
	\item Algunas decisiones de la secuencia, como los umbrales temporales, el número de reintentos, la activación de los evaluadores y la elección de modelos por etapa, se basan en la literatura o en criterios de diseño. Experimentos que modificaran un solo factor permitirían estudiar su efecto por separado.
	\item La taxonomía de estados y los roles se derivaron de la literatura y se concretaron desde una perspectiva de ingeniería. La revisión por especialistas y la comparación con anotaciones expertas permitirían valorar esta formulación y examinar si una etiqueta única representa adecuadamente discusiones que combinan varios estados.
	\item La revisión utilizó un único LLM de propósito general como evaluador. Las 46 intervenciones evaluadas ofrecen una primera referencia para repetir la revisión con otros modelos, contrastar sus puntuaciones con evaluaciones humanas y estudiar la calibración de la confianza y la reproducibilidad de las salidas.
	\item La independencia de la plataforma se examinó con archivos locales y una prueba de concepto con Open edX. Una segunda integración y un flujo de revisión y publicación supervisado por el instructor permitirían ampliar esta comprobación. La función de publicación actual se probó con componentes simulados.
	\item Los experimentos no incluyeron publicaciones con instrucciones dirigidas al modelo. Además, no se calculó el coste completo de una intervención ni se estudió la distribución de su latencia. Estos aspectos pueden incorporarse a futuras evaluaciones con más casos y configuraciones.
\end{itemize}

\section{Trabajo futuro}

La evaluación de las 46 intervenciones debe repetirse con varios modelos evaluadores, manteniendo la misma rúbrica y los mismos casos. La comparación debe examinar la coincidencia de las puntuaciones y las alertas, así como la estabilidad de sus justificaciones. De este modo podrá estudiarse cuánto dependen los resultados del modelo elegido como evaluador.

La elección del modelo debe estudiarse por separado en cada etapa. Cada experimento modificaría solamente el modelo utilizado en la clasificación, la decisión de intervención, la selección del rol o la generación de la respuesta. Los hilos de discusión, las instrucciones, el contexto y los modelos de las demás etapas se mantendrían fijos. Así podrían compararse la finalización de la etapa, el cumplimiento del esquema y la evaluación de la decisión o respuesta producida.

Los cambios en el contexto requieren experimentos separados. Cada prueba debería modificar un único elemento del contexto y mantener fija la configuración restante para poder estudiar su efecto.

La evaluación también debe incluir hilos de discusión en español y contextos educativos con convenciones de participación diferentes. Para ello se necesitan conjuntos de hilos de discusión y revisores familiarizados con cada contexto. El análisis debería comprobar si los estados, la decisión de intervenir, la selección del rol y la técnica y la respuesta propuesta siguen siendo adecuados.

También sería necesario ampliar la evaluación con datos reales procedentes de plataformas educativas y con discusiones asíncronas más actuales. Esta ampliación permitiría comprobar el comportamiento del sistema en condiciones de uso distintas de las representadas por los hilos empleados en este trabajo.

Las ejecuciones realizadas aportan evidencia de la viabilidad técnica de una arquitectura modular para facilitar discusiones académicas asíncronas. Los experimentos usaron hilos versionados y modelos locales y externos. Por separado, la prueba de concepto con Open edX adaptó hilos al mismo contrato de entrada. Queda por evaluar si este apoyo mejora la experiencia de los estudiantes y permite al instructor atender más discusiones sin reducir la calidad de la enseñanza.
